% Part: incompleteness
% Chapter: theories-computability
% Section: extensions-of-q-not-decidable

\documentclass[../../../include/open-logic-section]{subfiles}

\begin{document}

\olfileid{inc}{tcp}{cqn}

\olsection{$\Th{Q}$ دیاں بے تضاد توسیعاں ناقابلِ فیصلہ نیں}

\begin{explain}
یاد رکھو کہ اک نظریہ \emph{بے تضاد} اے جے اوہ
$!A$ تے $\lnot !A$ اکٹھے ثابت نہ کرے، کسے وی فارمولے~$!A$ لئی۔
تضاد توں ہر شے اخذ ہو سکدی اے۔ ایس لئی تضاد والے نظریے وچ ہر
جملہ ثابت ہو سکدا اے۔ ظاہر اے کہ جے نظریہ $\omega$-بے تضاد
ہووے، تاں اوہ بے تضاد وی اے۔ عام بے تضادی دی شرط کمزور اے۔
یعنی کجھ نظریے بے تضاد ہوندے نیں پر $\omega$-بے تضاد نہیں
ہوندے۔ اسیں \olref[oqn]{thm:oconsis-q} وچ مفروضے دی جگہ
عام بے تضادی رکھ کے ہور مضبوط قضیہ حاصل کر سکدے آں۔
\end{explain}

\begin{lem}
کوئی “ہمہ گیر قابلِ حساب تعلق” موجود نہیں۔ یعنی کوئی دو رکنی
قابلِ حساب تعلق $R(x,y)$ موجود نہیں جیہدے وچ ایہ خاصیت ہووے۔
جدوں وی $S(y)$ اک رکنی قابلِ حساب تعلق ہووے، کوئی ایہو جیہا $k$
ہووے جیہدے لئی ہر $y$ اُتے، $S(y)$ اُدوں تے صرف اُدوں سچ
ہووے جدوں $R(k,y)$ سچ ہووے۔
\end{lem}

\begin{proof}
فرض کرو $R(x,y)$ ہمہ گیر قابلِ حساب تعلق اے۔ $S(y)$ نوں
تعلق $\lnot R(y,y)$ منّ لئیے۔ کیوں جے $S(y)$ قابلِ حساب اے،
ایس لئی کسے $k$ لئی $S(y)$ تے $R(k,y)$ ہم ارز نیں۔
فیر $S(k)$، $R(k,k)$ تے $\lnot R(k,k)$ دونہاں دے ہم ارز
اے، جیہڑا تضاد اے۔
\end{proof}

\begin{thm}
فرض کرو $\Th{T}$ کوئی وی بے تضاد نظریہ اے جیہڑا $\Th{Q}$
نوں شامل کردا اے۔ تاں $\Th{T}$ قابلِ فیصلہ نہیں۔
\end{thm}

\begin{proof}
فرض کرو $\Th{T}$، $\Th{Q}$ دی اک بے تضاد تے قابلِ فیصلہ
توسیع اے۔ اسیں $\Th{T}$ دی مدد نال اک ہمہ گیر قابلِ حساب
تعلق متعین کراں گے۔ ایس توں تضاد نکلے گا۔

تعلق $R(x,y)$ اُدوں تے صرف اُدوں قائم اے جدوں
\begin{quote}
$x$، فارمولے $!D(u)$ دا کوڈ اے تے $\Th{T}$،
فارمولا $!D(\num y)$ ثابت کردا اے۔
\end{quote}
فرض مطابق $\Th{T}$ قابلِ فیصلہ اے، ایس لئی $R$ قابلِ حساب اے۔
ہن وکھائیے کہ $R$ ہمہ گیر اے۔ جے $S(y)$ کوئی وی قابلِ حساب
تعلق اے، تاں $\Th{Q}$ وچ، تے ایس لئی $\Th{T}$ وچ وی، اوہدی
نمائندگی فارمولا $!D_S(u)$ کر سکدا اے۔ فیر ہر $n$ لئی
\begin{eqnarray*}
S(\num n) & \lif & T \vdash !D_S(\num n) \\
& \lif & R(\Gn{!D_S(u)}, n)
\end{eqnarray*}
تے
\begin{eqnarray*}
\lnot S(\num n) & \lif & T \vdash \lnot !D_S(\num n) \\
& \lif & T \not\vdash !D_S(\num n) \quad \text{(کیوں جے $\Th{T}$
  بے تضاد اے)} \\
& \lif & \lnot R(\Gn{!D_S(u)}, n).
\end{eqnarray*}
یعنی ہر $y$ لئی $S(y)$ اُدوں تے صرف اُدوں سچ اے جدوں
$R(\Gn{!D_S(u)}, y)$ سچ اے۔ ایس لئی $R$ ہمہ گیر اے،
تے ایہی اوہ تضاد اے جیہڑا اسیں لبھ رہے ساں۔
\end{proof}

“سچا حساب” اوہ نظریہ اے جیہڑا $\Setabs{!A}{ \Sat{\Nat}{!A}}$
نال دتا جاندا اے، یعنی حساب دی زبان دے اوہناں جملیاں دا سیٹ
جیہڑے معیاری تعبیر وچ سچے نیں۔

\begin{cor}
سچا حساب قابلِ فیصلہ نہیں۔
\end{cor}

%حصے~\olref{undefinability:truth:section} وچ اسیں تارسکی دا
%ہور مضبوط نتیجہ بیان کراں گے۔

\end{document}
